\answer{ex:20:1}{$t_{\text{thức}}=\tau_r\ln\frac{1-L}{1-H}=16.8$~giờ, $t_{\text{ngủ}}=\tau_d\ln\frac HL=5.8$~giờ, chu kỳ $22.5$~giờ; dao động ngày đêm của các ngưỡng kéo bộ dao động về $24$~giờ, đó là khóa pha~\eqref{eq:pll}.}
\answer{ex:20:2}{(a)~$L=0.111$; với $L=0.092$ giấc ngủ sẽ kéo dài $9.6$~giờ. (b)~Tăng $22\%$ ($1/0.82=1.22$), chứ không phải $18\%$.}
\answer{ex:20:3}{$H=\frac{p-1}{p-1/q}=0.623$, $L=H/q=0.093$, trong đó $p=e^{16/\tau_r}$, $q=e^{8/\tau_d}$. Sau khi bị đánh thức sau $4$~giờ, pha bị dịch vĩnh viễn (ngủ sớm hơn $7.2$~giờ); với ngưỡng dao động, pha trở lại sau hai đến ba ngày.}
\answer{ex:20:4}{(a)~$r_\pm=\frac12\bigl(1\pm\sqrt{1-4k/w}\bigr)=0.947,\ 0.053$; $a_{\text{trên}}=3.51$, $a_{\text{dưới}}=0.49$. (b)~Làm việc $\approx0.33$~s, im lặng $\approx0.59$~s, chu kỳ $\approx0.93$~s, tỉ lệ làm việc $0.36$.}
\answer{ex:20:5}{$a_{\text{trên}}/r_+<b<a_{\text{dưới}}/r_-$: $3.71<b<9.20$. \nt{ACET} làm $b$ giảm xuống dưới $3.71$: giao điểm chuyển lên nhánh trên, và trạng thái làm việc ($r\approx1$) trở nên ổn định.}
\answer{ex:20:6}{$4.9$, $6.8$, $8.8$, $5.5$, $8.9$~giờ với $D=24$, $32$, $40$, $48$, $64$: độ dài do pha ngày đêm lúc đi ngủ quyết định, chứ không phải do món nợ.}
\answer{ex:20:7}{Sau B, sai số trên A trung bình là $0.51$ (từ $0.19$ đến $1.08$ qua $20$ bộ dữ liệu; đáp ứng bằng không cho $1$). Khi trộn lẫn, cả hai sai số đều nhỏ hơn $10^{-4}$.}
\answer{ex:20:8}{Bài mở. Co giãn đều giữ nguyên tỉ lệ giữa các trọng số, nhưng chỉ giữ nguyên đáp ứng của nơ-ron có ngưỡng khi ngưỡng cũng được co giãn cùng tỉ lệ; phát lại trộn lẫn thì làm thay đổi tỉ lệ. Việc các tiếp xúc lớn không đổi mâu thuẫn với sự co giãn hoàn toàn đều.}
